\section{Configuración, compilación y comprobación}

Estos comandos usan el proyecto existente. Son instrucciones de reproducción; no se ejecutaron nuevas mediciones durante la elaboración del informe.

\subsection{Comprobar el software y los nodos}
\begin{lstlisting}[language=bash]
cd /home/alumno16/experimentos_96
mpirun --version
mpicc --showme:command
lscpu

for h in 10.7.50.203 10.7.50.201 10.7.50.204; do
  ssh -F ssh_config "$h" \
    'hostname; nproc; mpirun --version | head -1'
done

# Diagnostico de conectividad; ping no mide ancho de banda.
ping -c 3 10.7.50.203
ping -c 3 10.7.134.58
\end{lstlisting}

\subsection{Compilar los dos programas}
\begin{lstlisting}[language=bash]
cd /home/alumno16/experimentos_96
mkdir -p /var/tmp/mpi_trapecio /var/tmp/mpi_bloques

mpicc -O3 -march=native -funroll-loops \
  src/trapecio.c -o /var/tmp/mpi_trapecio/trapecio -lm

mpicc -O3 -march=native -funroll-loops \
  src/mpi_matrix_2d.c \
  -o /var/tmp/mpi_bloques/mpi_matrix_2d -lm
\end{lstlisting}
\texttt{-O3} activa optimizaciones; \texttt{-march=native} permite instrucciones de la CPU local. El despliegue presupone los equipos compatibles documentados. \archivo{/var/tmp} conserva normalmente los archivos entre reinicios, aunque una política de limpieza puede borrarlos; comprobar su existencia sigue siendo necesario.

\subsection{Copiar y verificar los binarios}
\begin{lstlisting}[language=bash]
for h in 10.7.50.203 10.7.50.201 10.7.50.204; do
  ssh -F ssh_config "$h" \
    'mkdir -p /var/tmp/mpi_trapecio /var/tmp/mpi_bloques'
  scp -F ssh_config /var/tmp/mpi_trapecio/trapecio \
    "$h:/var/tmp/mpi_trapecio/"
  scp -F ssh_config /var/tmp/mpi_bloques/mpi_matrix_2d \
    "$h:/var/tmp/mpi_bloques/"
done
sha256sum /var/tmp/mpi_trapecio/trapecio \
          /var/tmp/mpi_bloques/mpi_matrix_2d
for h in 10.7.50.203 10.7.50.201 10.7.50.204; do
  ssh -F ssh_config "$h" \
    'sha256sum /var/tmp/mpi_trapecio/trapecio /var/tmp/mpi_bloques/mpi_matrix_2d'
done
\end{lstlisting}

\section{Comandos de lanzamiento y medición}

\subsection{Una ejecución, paso a paso}
La interfaz de los dos lanzadores es \texttt{experimento tamaño procesos}. Se aceptan los procesos 1, 24, 48, 72 y 96.
\begin{lstlisting}[language=bash]
cd /home/alumno16/experimentos_96

# Referencia de un proceso, nodo completo y cuatro nodos.
./lanzar.sh trapecio 100000000000 1
./lanzar.sh trapecio 100000000000 24
./lanzar.sh trapecio 100000000000 96

# Matrices: misma version 1d tiled 64 en todos los casos.
./lanzar.sh matrices 3072 1
./lanzar.sh matrices 3072 24
./lanzar.sh matrices 3072 96
./lanzar_wifi.sh matrices 3072 96
\end{lstlisting}
Para ver el comando expandido y la afinidad basta conservar la salida del lanzador. Las líneas \texttt{COMANDO} de los logs documentan qué se ejecutó; no basta el nombre de la carpeta.

\subsection{Batería completa y CSV}
\begin{lstlisting}[language=bash]
# Mostrar el plan; no ejecuta los experimentos.
./medir.sh todos --plan
./medir_wifi.sh todos --plan

# Ejecutar una red por vez, sin otra bateria simultanea.
./medir.sh todos
./medir_wifi.sh todos

# Tambien se puede elegir solo un experimento.
./medir.sh trapecio
./medir_wifi.sh matrices
\end{lstlisting}
Cada batería crea un directorio nuevo con fecha y sufijo aleatorio, un log por prueba y \archivo{tiempos.csv}. El CSV contiene \texttt{experimento,tamano,procesos,T\_Total}. Una fila es una ejecución; no se hacen promedios ocultos. Al terminar se invoca \archivo{graficar.py}. Si una ejecución falla o no entrega un tiempo válido, el script se detiene.

\subsection{Qué hace el script de medición}
El script recorre tamaños y procesos, ejecuta el lanzador, extrae \texttt{T\_Total} de la salida y lo añade al CSV. El cronómetro está en C; \texttt{sed} solamente lee su resultado. No hace falta \texttt{time} del shell para obtener el tiempo del algoritmo.

\subsection{Regenerar gráficos de las corridas actuales}
Los siguientes comandos reutilizan datos ya guardados:
\begin{lstlisting}[language=bash]
cd /home/alumno16/experimentos_96
eth=resultados/20260924_202548_RnIhhM/tiempos.csv
wifi=resultados_wifi/20260924_203957_VSNY0P/tiempos.csv

python3 graficar.py "$eth"
python3 graficar.py "$wifi" --red Wi-Fi
python3 comparar_eth_vs_wifi.py "$eth" "$wifi" \
  --salida comparacion_eth_vs_wifi
\end{lstlisting}
Se generan tablas de métricas y gráficos de tiempo, speedup y eficiencia. La comparación calcula $T_W/T_E$. Los gráficos de este informe se regeneran con un script propio que toma copias de los mismos datos y de los históricos, con estilos aptos para impresión en blanco y negro.

\subsection{Reconstruir este informe}
\begin{lstlisting}[language=bash]
cd /home/alumno16/cluster-mpi-96-optimizaciones/revision
./compilar.sh
\end{lstlisting}
\archivo{compilar.sh} ejecuta tres scripts y compila \archivo{informe_mejorado.tex}: \archivo{generar_datos.py} (tablas y figuras de las baterías y de las series históricas), \archivo{generar_optimizacion.py} (tablas y figuras de la campaña) y \archivo{auditar_evidencia.py} (auditoría y variabilidad). Los tres leen la evidencia de la carpeta superior del mismo repositorio sin modificarla; otra ubicación se indica con la variable \texttt{EVIDENCIA\_MPI}. No se llama a \texttt{mpirun}, no se cambian las configuraciones de los workers ni se repiten los experimentos.

\subsection{Lectura de un resultado real}
La línea siguiente se toma del log actual de matrices Wi-Fi, $N=3072$, $P=96$. Se muestran solo los campos necesarios para leer las fases:
\begin{lstlisting}
N=3072, Procs=96, Algo=1d, Kernel=tiled, BS=64
T_Dist=95.931070, T_Calc=0.100977
T_Gather=58.381650, T_Total=101.042337
\end{lstlisting}
\texttt{T\_Dist} reúne \texttt{Scatterv} y \texttt{Bcast}; la versión usada en las baterías no registra cada una por separado. La versión de la campaña añade \texttt{T\_Scatter} y \texttt{T\_Bcast}. El valor de \texttt{T\_Gather} puede incluir espera mientras otros procesos aún distribuyen o calculan. Por eso no se suma a los otros máximos para reconstruir el total.

\section{Cómo registrar los MB en una próxima ejecución}

La primera batería no guarda MB; la segunda registra solo TX+RX del servidor. Para una futura prueba puntual que mida las cuatro interfaces se incluye \archivo{medir_trafico.sh}. Lee los contadores de la interfaz elegida en los cuatro nodos y llama al lanzador existente una sola vez. Consulta los workers por Ethernet para reducir el tráfico de control agregado a una medición Wi-Fi; al medir Ethernet esas consultas sí pueden contribuir al contador. Sigue siendo una medida de interfaces, sin aislamiento de tráfico ajeno.

\begin{lstlisting}[language=bash]
cd /home/alumno16/cluster-mpi-96-optimizaciones

# Ejemplo: UNA ejecucion adicional con registro de MB.
# No forma parte de los resultados de este informe.
./medir_trafico.sh wifi matrices 3072 96

# Para Ethernet, cambiar solo el primer argumento:
./medir_trafico.sh ethernet matrices 3072 96
\end{lstlisting}
El nuevo script deja \archivo{trafico.csv}, \archivo{antes.tsv}, \archivo{despues.tsv} y \archivo{ejecucion.log} en una carpeta nueva. Los archivos de antes y después conservan los contadores brutos, además del total calculado. No afirma que sus MB equivalgan solo a la carga útil de MPI ni a los bytes radiados.

\subsection{Aislar el efecto de una opción}
Para estudiar específicamente \texttt{vader} o el algoritmo de \texttt{Bcast}, no se debe comparar 16 procesos \texttt{ikj} con 96 \texttt{tiled} y atribuir toda la diferencia a una opción. Haría falta conservar el mismo tamaño, kernel, número de procesos y asignación. Los comandos de consulta siguientes permiten inventariar los parámetros disponibles en la versión instalada:
\begin{lstlisting}[language=bash]
ompi_info --version
ompi_info --param btl vader --level 9
ompi_info --param coll tuned --level 9
\end{lstlisting}
Consultar parámetros no demuestra cuál se eligió en una corrida anterior. Habría que registrar también la selección durante la ejecución o forzar explícitamente un algoritmo compatible en un experimento controlado. Es una comprobación pendiente.

\section{Scripts actuales conservados}

Las copias de esta sección corresponden a los lanzadores y a la medición de la batería analizada. Los archivos de configuración SSH y los demás scripts acompañan al informe en \archivo{scripts/}; no se incluyen contraseñas ni claves privadas.

\subsection{Lanzador Wi-Fi}
El lanzador Ethernet usa el mismo esquema de asignación con las IP de cable. Se incluye el Wi-Fi completo para dejar visibles las opciones que interesan en la comparación histórica.
\lstinputlisting[language=bash,basicstyle=\ttfamily\scriptsize]{scripts/lanzar_wifi.sh}

\subsection{Medición Bash y generación de gráficos}
\lstinputlisting[language=bash,basicstyle=\ttfamily\scriptsize]{scripts/medir.sh}
La variante \archivo{medir_wifi.sh} llama a \archivo{lanzar_wifi.sh}, guarda en \archivo{resultados_wifi/} e identifica la red al generar los gráficos. Ambos scripts validan que el tiempo exista y sea positivo; \textbf{no validan automáticamente los checksums}.

\section{Archivos de respaldo y fuentes}

El repositorio \archivo{cluster-mpi-96-optimizaciones} incluye una copia de la evidencia necesaria para reproducir las tablas y los gráficos; las rutas de la tabla son relativas a la carpeta principal. Esta versión revisada no duplica la evidencia: la lee desde la carpeta superior. \archivo{datos/procedencia.json} relaciona cada copia con su ruta original y su hash SHA-256; \archivo{datos/validacion.json} resume las comprobaciones. \archivo{datos/metricas_informe.csv} conserva las métricas calculadas sin el redondeo de las tablas.

\begin{center}\small
\begin{tabularx}{\linewidth}{>{\raggedright\arraybackslash}p{5.5cm}X}
\toprule
Ruta en el repositorio & Contenido \\
\midrule
\archivo{datos/ethernet/} & 25 resultados de la primera batería y sus logs por cable. \\
\archivo{datos/wifi/} & 25 resultados de la primera batería y sus logs Wi-Fi. \\
\archivo{datos/ampliada/} & 75 logs y cuatro archivos de resultados de la segunda batería. \\
\archivo{datos/procedencia_ampliada.json} & Rutas y SHA-256 de los 79 archivos nuevos de evidencia y tres fuentes añadidas. \\
\archivo{datos/historicos/resultados_profiling_repeticion.json} & Primera serie de madrugada; contiene los 906.351 s con 8 procesos. \\
\archivo{datos/historicos/resultados_bloques.json} & Caché, escalado local y serie de madrugada con MB y dos repeticiones en el clúster. \\
\archivo{datos/historicos/resultados_matrices_eth_vs_wifi.json} & Serie del mediodía con tiempos y MB por red. \\
\archivo{src/} & Fuentes de trapecio, matrices y kernels. \archivo{mpi_matrix_2d.c} es la versión de la campaña (con \texttt{1d-hier}); la usada en las baterías está en \archivo{experimentos_96/src/} y en el commit \texttt{f81f488}. \\
\archivo{resultados_optimizacion/} & CSV, logs, \archivo{entorno.json} y \archivo{SHA256SUMS} de la campaña de optimización. \\
\archivo{fuentes/INFORME_ANALISIS_EXPANDIDO_96.md} & Borrador de análisis conservado como fuente, con sus afirmaciones revisadas. \\
\archivo{fuentes/exploratorios/} & Prototipos de núcleos P/E, sin salidas individuales verificables en la batería. \\
\archivo{scripts/} & Copias de lanzadores, medición, gráficos y configuración SSH. \\
\archivo{generar_datos.py} & Lee la evidencia, valida consistencia y genera tablas y seis figuras. Esta versión usa una copia adaptada que solo escribe en su propia carpeta. \\
\archivo{medir_trafico.sh} & Herramienta opcional para una futura medición de MB. \\
\bottomrule
\end{tabularx}\end{center}

El contexto actualizado del proyecto orientó la revisión, pero las cifras finales se tomaron de los CSV y logs. Se corrigieron diferencias de tiempo y atribuciones entre series: 1.053217 s para trapecio grande Wi-Fi de la primera batería, 906.351 s en la serie original con 8 procesos, y 1112 MB en la serie posterior con 16 procesos. La primera batería no midió MB; la segunda midió TX+RX del servidor, que no equivale al TX agregado histórico.

\begin{thebibliography}{9}
\raggedright
\bibitem{foster} Ian Foster. \emph{Designing and Building Parallel Programs}, 1995. Capítulo 2: Designing Parallel Algorithms. \url{https://www.mcs.anl.gov/~itf/dbpp/text/node14.html}.
\bibitem{ompi-tcp} Open MPI. \emph{FAQ: Tuning the run-time characteristics of MPI TCP communications}. Documentación para Open MPI 4.x y anteriores; apartados sobre transportes locales y selección de interfaces. \url{https://www.open-mpi.org/faq/?category=tcp}.
\bibitem{ompi-source} Open MPI, versión 4.1.6. \emph{coll\_tuned\_decision\_fixed.c}, función \texttt{ompi\_coll\_tuned\_bcast\_intra\_dec\_fixed}. Código fuente oficial. \url{https://github.com/open-mpi/ompi/blob/v4.1.6/ompi/mca/coll/tuned/coll_tuned_decision_fixed.c}.
\bibitem{datos} Registros del laboratorio del 24 de septiembre de 2026. Archivos CSV, JSON, logs y fuentes enumerados en este anexo; copias y hashes incluidos en el paquete del informe.
\end{thebibliography}
